\section{Conclusión}

Se desarrolló el diseño sísmico de los muros Pier~C y Pier~J, y de las vigas sísmicas B7, B8 y B39
del edificio, conforme a ACI~318-19, DS~60 y NCh433.

En cuanto a las compresiones en los pisos inferiores, ambos muros presentan tensiones de compresión
media que se mantienen dentro del límite admisible de $0{,}35 f_c'$. El valor más exigente corresponde
al Pier~C en piso~1, con $\sigma_u = 70{,}7\;\text{kg/cm}^2$, que representa el 57,7\,\% del límite.
Los niveles superiores presentan compresiones significativamente menores, lo que refleja la reducción
progresiva de la carga axial hacia el techo.

Respecto a las armaduras de corte, el diseño con sobrerresistencia $\Omega_v = 2{,}0$ incrementa
considerablemente la demanda de cortante de diseño respecto al valor de la envolvente, lo que resulta
en cuantías transversales superiores al mínimo normativo en los pisos bajos de ambos muros. En
Pier~C la cuantía máxima se alcanza en piso~2 ($\rho_t = 0{,}0045$), mientras que en Pier~J ocurre
en piso~2 ($\rho_t = 0{,}0057$). A partir del piso~5 en adelante la armadura mínima de
$\rho_t = 0{,}0025$ gobierna en ambos muros, correspondiente a malla doble $\phi10$ @ $250\;\text{mm}$
por cara.

En lo que respecta a las armaduras de flexo-compresión, el método de Cárdenas-Magura entrega
cuantías verticales de $\rho_v = 0{,}0035$ en todos los pisos analizados para ambos muros, lo que
corresponde a $A_{s,w} = 47{,}25\;\text{cm}^2$ en Pier~C y $A_{s,w} = 50{,}40\;\text{cm}^2$ en
Pier~J. La profundidad del eje neutro se mantiene en todos los casos por debajo de $3l_w/8$,
confirmando comportamiento dúctil de la sección.

En relación a la calidad del hormigón, el uso de G35 en los pisos 1 al~3 resulta determinante para
satisfacer las verificaciones de compresión máxima y corte máximo admisible en la zona más solicitada
del muro. En los pisos superiores, el hormigón G25 es suficiente dado el nivel de solicitaciones.

Respecto a la posible incursión inelástica, la verificación de confinamiento indica que el Pier~C en
piso~1 supera el umbral de tensiones $\sigma > 0{,}2f_c'$, lo que exige la disposición de elementos
especiales de borde confinados con $l_\text{conf} = 31{,}4\;\text{cm}$ y separación máxima
$s_\text{máx} = 7{,}2\;\text{cm}$, extendidos verticalmente una altura de $5{,}40\;\text{m}$ desde
la base. El resto de los pisos de Pier~C y la totalidad de Pier~J no requieren confinamiento. Las
deformaciones máximas de compresión obtenidas son en todos los casos inferiores a $0{,}003$, lejos
del límite de $0{,}008$ establecido por el DS~60, lo que confirma que la incursión inelástica es
acotada y controlada.

En cuanto al diseño de vigas sísmicas, las tres vigas cumplen los requisitos de resistencia a flexión
y corte en todos los pisos. La viga B39 en piso~1 es el elemento más solicitado de la estructura,
con $V_u = 77{,}79\;\text{tonf}$ y $M_u = 29{,}21\;\text{tonf}{\cdot}\text{m}$, y requiere la
armadura más pesada del conjunto: $3\phi25$ longitudinal con estribos $\phi12$ @ $0{,}30\;\text{m}$.
Las vigas B7 y B8 presentan solicitaciones notablemente menores, gobernadas por la armadura mínima
en los pisos superiores.